%=========== ΑΠΑΝΤΗΣΕΙΣ - 250 SMART ΣΩΣΤΟ ΛΑΘΟΣ ===========
\section{Απαντήσεις στις Ερωτήσεις Σωστό - Λάθος}

\subsection{Όρια συνέχεια}

\begin{enumerate}[label=\textbf{1.\arabic*.}]
\item \textbf{Λάθος}. Πρέπει να ισχύει $x^2>0$, άρα $x\neq0$. Επομένως το πεδίο ορισμού είναι
\[
\mathbb{R}^*=(-\infty,0)\cup(0,+\infty).
\]

\item \textbf{Λάθος}. Ισχύει
\[
\sqrt{x^2}=|x|.
\]
Άρα για $x<0$ έχουμε $\sqrt{x^2}=-x\neq x$.

\item \textbf{Λάθος}. Το ίδιο πεδίο ορισμού και οι ίδιες ακραίες τιμές δεν καθορίζουν μοναδικά μια συνάρτηση. Για παράδειγμα, στο $[0,1]$ οι συναρτήσεις $f(x)=x$ και $g(x)=1-x$ έχουν ίδια ελάχιστη και μέγιστη τιμή, αλλά δεν είναι ίσες.

\item \textbf{Σωστό}. Το πεδίο ορισμού της $f+g$ είναι η τομή των πεδίων ορισμού των $f$ και $g$. Αν αυτή η τομή είναι όλο το $\mathbb{R}$, τότε και οι δύο συναρτήσεις ορίζονται σε όλο το $\mathbb{R}$.

\item \textbf{Λάθος}. Το πεδίο ορισμού της $f+g$ είναι γενικά η τομή των πεδίων ορισμού των $f$ και $g$, όχι η ένωση.

\item \textbf{Λάθος}. Η
\[
f(x)=\frac{x^2-1}{x-1}
\]
έχει πεδίο ορισμού το $\mathbb{R}\setminus\{1\}$, ενώ η $g(x)=x+1$ έχει πεδίο ορισμού το $\mathbb{R}$. Άρα δεν ταυτίζονται ως συναρτήσεις.

\item \textbf{Σωστό}. Για $x>0$ έχουμε
\[
\ln x^2=\ln(x^2)=2\ln x.
\]

\item \textbf{Λάθος}. Για $x<0$ το $\ln x$ δεν ορίζεται στους πραγματικούς αριθμούς. Η σωστή σχέση για $x\neq0$ είναι
\[
\ln(x^2)=2\ln|x|.
\]

\item \textbf{Σωστό}. Αν $f,g$ είναι άρτιες, τότε
\[
(fg)(-x)=f(-x)g(-x)=f(x)g(x)=(fg)(x).
\]
Άρα το γινόμενό τους είναι άρτια συνάρτηση.

\item \textbf{Σωστό}. Αν η $f$ είναι περιττή και $0$ ανήκει στο πεδίο ορισμού της, τότε
\[
f(0)=-f(0),
\]
οπότε
\[
f(0)=0.
\]

\item \textbf{Λάθος}. Μια συνάρτηση αντιστοιχίζει σε κάθε $x$ το πολύ μία τιμή. Ο άξονας $y'y$ έχει $x=0$, άρα η γραφική παράσταση μπορεί να τον τέμνει το πολύ σε ένα σημείο.

\item \textbf{Σωστό}. Αν η $f$ είναι άρτια, τότε $f(x)=f(-x)$. Αν το πεδίο ορισμού περιέχει δύο διαφορετικά συμμετρικά σημεία $x$ και $-x$, η $f$ δεν είναι $1-1$. Εξαίρεση υπάρχει μόνο όταν δεν υπάρχουν δύο τέτοια διαφορετικά σημεία.

\item \textbf{Σωστό}. Αν η $f$ έχει μη μηδενική περίοδο $T$, τότε
\[
f(x+T)=f(x)
\]
για κατάλληλα διαφορετικά σημεία του πεδίου ορισμού της. Άρα δεν είναι $1-1$.

\item \textbf{Λάθος}. Από τη σχέση $f^2(x)=x^2$ προκύπτει μόνο ότι
\[
f(x)=x \quad \text{ή} \quad f(x)=-x
\]
για κάθε συγκεκριμένο $x$. Το πρόσημο μπορεί να αλλάζει από σημείο σε σημείο, π.χ. $f(x)=|x|$.

\item \textbf{Λάθος}. Δεν αρκεί η $f$ να είναι γνησίως αύξουσα. Αν η $f$ παίρνει και αρνητικές και θετικές τιμές χωρίς να μηδενίζεται στο πεδίο της, η $\frac1f$ δεν είναι υποχρεωτικά γνησίως φθίνουσα σε όλο το διάστημα.

\item \textbf{Σωστό}. Αν $x_1<x_2$ και η $f$ είναι γνησίως φθίνουσα, τότε
\[
f(x_1)>f(x_2).
\]
Επειδή και η $g$ είναι γνησίως φθίνουσα, από την τελευταία ανισότητα παίρνουμε
\[
g(f(x_1))<g(f(x_2)).
\]
Άρα η $g\circ f$ είναι γνησίως αύξουσα.

\item \textbf{Σωστό}. Αν $g(x_1)=g(x_2)$, τότε
\[
(f\circ g)(x_1)=(f\circ g)(x_2).
\]
Επειδή η $f\circ g$ είναι $1-1$, προκύπτει $x_1=x_2$. Άρα η $g$ είναι $1-1$.

\item \textbf{Λάθος}. Η $f$ χρειάζεται να είναι $1-1$ μόνο πάνω στο σύνολο τιμών της $g$, όχι απαραίτητα σε όλο το πεδίο ορισμού της.

\item \textbf{Λάθος}. Το άθροισμα μιας γνησίως αύξουσας και μιας γνησίως φθίνουσας συνάρτησης δεν είναι αναγκαστικά σταθερό. Για παράδειγμα, $f(x)=2x$ και $g(x)=-x$ δίνουν $f(x)+g(x)=x$.

\item \textbf{Σωστό}. Η γραφική παράσταση της $-f(x)$ προκύπτει από την $C_f$ με αλλαγή προσήμου στις τεταγμένες. Άρα είναι συμμετρική της $C_f$ ως προς τον άξονα $x'x$.

\item \textbf{Λάθος}. Η γραφική παράσταση της $f(-x)$ είναι συμμετρική της $C_f$ ως προς τον άξονα $y'y$, όχι ως προς τον άξονα $x'x$.

\item \textbf{Σωστό}. Αν η $f$ έχει αντίστροφη, τότε είναι $1-1$. Άρα για κάθε $y$ του συνόλου τιμών της, η εξίσωση $f(x)=y$ έχει ακριβώς μία λύση.

\item \textbf{Λάθος}. Οι γραφικές παραστάσεις των $f$ και $f^{-1}$ είναι συμμετρικές ως προς την ευθεία $y=x$, αλλά μπορούν να τέμνονται και εκτός της $y=x$. Για παράδειγμα, η $f(x)=1-x$ είναι ίση με την αντίστροφή της.

\item \textbf{Λάθος}. Η $f(x)=\sqrt{x-2}$ έχει πεδίο ορισμού το $[2,+\infty)$ και δεν είναι συμμετρική ως προς τον άξονα $y'y$.

\item \textbf{Σωστό}. Από τον ορισμό της αντίστροφης, αν
\[
f(a)=\beta,
\]
τότε
\[
f^{-1}(\beta)=a.
\]

\item \textbf{Λάθος}. Μια συνάρτηση μπορεί να είναι $1-1$ χωρίς να είναι υποχρεωτικά γνησίως μονότονη σε όλο το πεδίο ορισμού της.

\item \textbf{Λάθος}. Για να υπάρχει το
\[
\lim_{x\to x_0}f(x),
\]
δεν είναι απαραίτητο να ορίζεται η $f$ στο $x_0$. Αρκεί να ορίζεται κοντά στο $x_0$.

\item \textbf{Λάθος}. Από την ανισότητα $f(x)<g(x)$ κοντά στο $x_0$ προκύπτει το πολύ
\[
\lim_{x\to x_0}f(x)\leq\lim_{x\to x_0}g(x),
\]
όταν τα όρια υπάρχουν. Μπορεί τα δύο όρια να είναι ίσα.

\item \textbf{Λάθος}. Αν
\[
\lim_{x\to x_0}f(x)>0,
\]
τότε η $f$ είναι θετική κοντά στο $x_0$, όχι υποχρεωτικά σε όλο το πεδίο ορισμού της.

\item \textbf{Λάθος}. Από το
\[
\lim_{x\to x_0}|f(x)|=|\ell|\neq0
\]
δεν προκύπτει απαραίτητα ότι η $f$ έχει όριο. Μπορεί οι τιμές της να εναλλάσσονται ανάμεσα σε θετικές και αρνητικές με ίδιο μέτρο.

\item \textbf{Σωστό}. Αν
\[
\lim_{x\to x_0}f(x)=+\infty,
\]
τότε, κοντά στο $x_0$, οι τιμές της $f$ γίνονται μεγαλύτερες από οποιονδήποτε θετικό αριθμό. Άρα τελικά είναι θετικές.

\item \textbf{Λάθος}. Ισχύει
\[
\left|\frac{\hm x}{x}\right|\leq \frac{1}{x}
\]
για $x>0$. Άρα
\[
\lim_{x\to+\infty}\frac{\hm x}{x}=0,
\]
όχι $1$.

\item \textbf{Λάθος}. Δεν αρκεί να υπάρχει
\[
\lim_{x\to x_0}f(x)\geq0.
\]
Πρέπει η παράσταση $\sqrt{f(x)}$ να ορίζεται κοντά στο $x_0$, δηλαδή να έχουμε $f(x)\geq0$ κοντά στο $x_0$.

\item \textbf{Λάθος}. Αν $f(x)<0$ κοντά στο $x_0$, τότε η $\sqrt{f(x)}$ δεν ορίζεται στους πραγματικούς αριθμούς κοντά στο $x_0$. Άρα δεν έχει πραγματικό όριο $0$.

\item \textbf{Σωστό}. Για $x\to0^+$ έχουμε
\[
\frac{1}{x^3}\to+\infty,
\]
ενώ για $x\to0^-$ έχουμε
\[
\frac{1}{x^3}\to-\infty.
\]
Άρα το δίπλευρο όριο δεν υπάρχει ούτε ως $+\infty$ ούτε ως $-\infty$.

\item \textbf{Σωστό}. Καθώς $x\to0$, έχουμε $|x|\to0^+$, άρα
\[
\frac{1}{|x|}\to+\infty.
\]

\item \textbf{Λάθος}. Από το
\[
\lim_{x\to x_0}f(x)=+\infty
\]
προκύπτει ότι η $f$ είναι θετική κοντά στο $x_0$, όχι υποχρεωτικά σε όλο το πεδίο ορισμού της.

\item \textbf{Σωστό}. Αν
\[
\lim_{x\to+\infty}f(x)=-\infty,
\]
τότε για αρκετά μεγάλα $x$ ισχύει, για παράδειγμα, $f(x)<0$. Άρα η $f$ παίρνει αρνητικές τιμές.

\item \textbf{Λάθος}. Η μορφή
\[
0\cdot(+\infty)
\]
είναι απροσδιοριστία. Ανάλογα με τις συναρτήσεις, το γινόμενο μπορεί να έχει διαφορετικά όρια.

\item \textbf{Λάθος}. Η μορφή
\[
(+\infty)+(+\infty)
\]
δεν είναι απροσδιοριστία. Το άθροισμα τείνει στο $+\infty$.

\item \textbf{Σωστό}. Η μορφή
\[
(+\infty)-(+\infty)
\]
είναι απροσδιοριστία, γιατί το αποτέλεσμα εξαρτάται από τον τρόπο με τον οποίο αποκλίνουν οι δύο ποσότητες.

\item \textbf{Σωστό}. Αν η $g$ είναι φραγμένη κοντά στο $x_0$, υπάρχει $M>0$ με
\[
|g(x)|\leq M.
\]
Τότε
\[
|f(x)g(x)|\leq M|f(x)|\to0,
\]
άρα
\[
\lim_{x\to x_0}f(x)g(x)=0.
\]

\item \textbf{Λάθος}. Για τη σύνθεση χρειάζεται προσοχή όταν $g(x)=u_0$ κοντά στο $x_0$. Αν η $f$ δεν είναι συνεχής στο $u_0$, μπορεί το $f(g(x))$ να μην τείνει στο $\ell$.

\item \textbf{Σωστό}. Η ευθεία $x=x_0$ είναι κατακόρυφη ασύμπτωτη όταν τουλάχιστον ένα από τα πλευρικά όρια της $f$ στο $x_0$ είναι $+\infty$ ή $-\infty$.

\item \textbf{Σωστό}. Μια συνάρτηση μπορεί να είναι ορισμένη στο $x_0$ και ταυτόχρονα να έχει κατακόρυφη ασύμπτωτη $x=x_0$, αν κάποιο πλευρικό όριο είναι άπειρο.

\item \textbf{Λάθος}. Κάθε πολυωνυμική συνάρτηση είναι συνεχής σε όλο το $\mathbb{R}$. Άρα δεν μπορεί να έχει κατακόρυφη ασύμπτωτη.

\item \textbf{Λάθος}. Αν μηδενίζεται ο παρονομαστής μιας ρητής συνάρτησης, μπορεί να υπάρχει αιρόμενη ασυνέχεια και όχι κατακόρυφη ασύμπτωτη. Για παράδειγμα,
\[
\frac{x^2-1}{x-1}=x+1,\quad x\neq1.
\]

\item \textbf{Σωστό}. Η ευθεία
\[
y=\lambda x+\beta
\]
είναι πλάγια ασύμπτωτη της $C_f$ στο $+\infty$, όταν
\[
\lim_{x\to+\infty}\bigl[f(x)-(\lambda x+\beta)\bigr]=0.
\]

\item \textbf{Λάθος}. Αν η $f$ είχε δύο διαφορετικές οριζόντιες ασύμπτωτες στο $+\infty$, τότε το ίδιο όριο στο $+\infty$ θα έπρεπε να ισούται με δύο διαφορετικούς αριθμούς, πράγμα αδύνατο.

\item \textbf{Λάθος}. Στο ίδιο άκρο, δηλαδή στο $+\infty$, μια συνάρτηση δεν μπορεί να έχει ταυτόχρονα οριζόντια και πλάγια ασύμπτωτη με μη μηδενική κλίση.

\item \textbf{Λάθος}. Αν υπήρχαν δύο διαφορετικές πλάγιες ασύμπτωτες στο $+\infty$, η διαφορά τους θα έπρεπε να τείνει στο $0$, κάτι που δεν γίνεται για δύο διαφορετικές ευθείες.

\item \textbf{Σωστό}. Είναι γνωστό ότι
\[
\lim_{x\to-\infty}e^x=0.
\]

\item \textbf{Σωστό}. Για $x\to0^+$ έχουμε
\[
\ln x\to-\infty.
\]

\item \textbf{Λάθος}. Μια συνάρτηση μπορεί να είναι συνεχής σε κάθε σημείο του πεδίου ορισμού της, χωρίς το πεδίο ορισμού της να είναι όλο το $\mathbb{R}$. Για παράδειγμα, η $\ln x$ είναι συνεχής στο $(0,+\infty)$.

\item \textbf{Λάθος}. Μια συνάρτηση μπορεί να είναι συνεχής σε άκρο του πεδίου ορισμού της, χωρίς να ορίζεται σε ολόκληρο ανοιχτό διάστημα γύρω από αυτό. Για παράδειγμα, η $\sqrt{x}$ είναι συνεχής στο $0$ στο πεδίο της.

\item \textbf{Λάθος}. Η ισότητα των πλευρικών ορίων εξασφαλίζει την ύπαρξη του ορίου, αλλά για συνέχεια χρειάζεται επιπλέον
\[
\lim_{x\to x_0}f(x)=f(x_0).
\]

\item \textbf{Σωστό}. Από τη συνέχεια στο $x_0$ και το $f(x_0)>0$, οι τιμές της $f$ παραμένουν θετικές σε μια περιοχή του $x_0$.

\item \textbf{Σωστό}. Η σύνθεση συνεχών συναρτήσεων είναι συνεχής στα σημεία όπου ορίζεται η σύνθεση.

\item \textbf{Λάθος}. Το άθροισμα δύο μη συνεχών συναρτήσεων μπορεί να είναι συνεχές. Για παράδειγμα, αν $g=-f$, τότε $f+g=0$.

\item \textbf{Λάθος}. Μπορεί η $|f|$ να είναι συνεχής ενώ η $f$ να μην είναι. Για παράδειγμα, αν $f(x)=1$ για $x\in\mathbb{Q}$ και $f(x)=-1$ για $x\notin\mathbb{Q}$, τότε $|f(x)|=1$.

\item \textbf{Σωστό}. Από το Θεώρημα Μέγιστης-Ελάχιστης Τιμής, κάθε συνεχής συνάρτηση σε κλειστό διάστημα $[a,\beta]$ παίρνει μέγιστη και ελάχιστη τιμή.

\item \textbf{Λάθος}. Μια συνεχής συνάρτηση σε ανοιχτό διάστημα μπορεί να παρουσιάζει μέγιστη και ελάχιστη τιμή. Για παράδειγμα, η $f(x)=0$ στο $(a,\beta)$ έχει και μέγιστη και ελάχιστη τιμή το $0$.

\item \textbf{Λάθος}. Το Θεώρημα Bolzano εξασφαλίζει τουλάχιστον μία ρίζα, όχι ακριβώς μία.

\item \textbf{Σωστό}. Αν η $f$ είναι συνεχής στο $[a,\beta]$ και δεν μηδενίζεται, τότε διατηρεί σταθερό πρόσημο στο διάστημα. Άρα τα $f(a)$ και $f(\beta)$ είναι ομόσημα και
\[
f(a)f(\beta)>0.
\]

\item \textbf{Λάθος}. Αν
\[
f(a)f(\beta)=0,
\]
τότε κάποια ρίζα βρίσκεται σε άκρο του διαστήματος. Το Θεώρημα Bolzano δεν εξασφαλίζει ρίζα στο εσωτερικό $(a,\beta)$.

\item \textbf{Σωστό}. Η εικόνα κλειστού διαστήματος μέσω συνεχούς συνάρτησης είναι κλειστό διάστημα
\[
[m,M],
\]
όπου $m$ και $M$ είναι η ελάχιστη και η μέγιστη τιμή της συνάρτησης.

\item \textbf{Σωστό}. Από το Θεώρημα Ενδιάμεσων Τιμών, μια συνεχής συνάρτηση σε κλειστό διάστημα παίρνει όλες τις τιμές ανάμεσα στην ελάχιστη και τη μέγιστη τιμή της.

\item \textbf{Σωστό}. Αν η $f$ είναι συνεχής και γνησίως αύξουσα στο $[a,\beta]$, τότε η ελάχιστη τιμή της είναι $f(a)$ και η μέγιστη $f(\beta)$. Άρα το σύνολο τιμών της είναι
\[
[f(a),f(\beta)].
\]

\item \textbf{Λάθος}. Η συνάρτηση
\[
f(x)=\frac{x^2-1}{x-1}
\]
δεν ορίζεται στο $x=1$. Άρα δεν είναι συνεχής στο $\mathbb{R}$.

\item \textbf{Λάθος}. Για $x\neq0$ έχουμε $f(x)=x$, οπότε
\[
\lim_{x\to0}f(x)=0.
\]
Όμως $f(0)=1$, άρα η συνάρτηση δεν είναι συνεχής στο $0$.

\item \textbf{Σωστό}. Κάθε πολυωνυμική συνάρτηση είναι συνεχής στο $\mathbb{R}$.

\item \textbf{Λάθος}. Μια ρίζα στο $(a,\beta)$ δεν συνεπάγεται αλλαγή προσήμου στα άκρα. Για παράδειγμα, η $f(x)=(x-\gamma)^2$ έχει ρίζα στο $\gamma\in(a,\beta)$, αλλά είναι μη αρνητική.

\item \textbf{Λάθος}. Αν δεν ισχύει η συνέχεια, δεν μπορούμε να εφαρμόσουμε Bolzano, αλλά δεν αποκλείεται να υπάρχει ρίζα στο $(a,\beta)$.

\item \textbf{Σωστό}. Κάθε συνεχής και $1-1$ συνάρτηση σε κλειστό διάστημα είναι υποχρεωτικά γνησίως μονότονη.

\item \textbf{Σωστό}. Μια συνάρτηση μπορεί να αλλάζει πρόσημο εκατέρωθεν μιας ρίζας χωρίς να είναι συνεχής εκεί. Για παράδειγμα,
\[
f(x)=
\begin{cases}
-1,&x<0,\\
0,&x=0,\\
1,&x>0.
\end{cases}
\]

\item \textbf{Λάθος}. Για το Θεώρημα Μέγιστης-Ελάχιστης Τιμής χρειάζεται η συνάρτηση να είναι συνεχής σε κλειστό διάστημα $[a,\beta]$, όχι απλώς σε ανοιχτό διάστημα.

\item \textbf{Σωστό}. Μια συνεχής συνάρτηση με πεδίο ορισμού ανοιχτό διάστημα μπορεί να έχει κλειστό σύνολο τιμών. Για παράδειγμα, η $f(x)=\hm x$ στο $(0,2\pi)$ έχει σύνολο τιμών το $[-1,1]$.

\item \textbf{Λάθος}. Η μορφή
\[
\frac{0}{0}
\]
είναι απροσδιοριστία. Για παράδειγμα, για $f(x)=x$ και $g(x)=2x$ στο $0$, το όριο του πηλίκου είναι $\frac12$, όχι $1$.

\item \textbf{Σωστό}. Αν η $f$ είναι γνησίως μονότονη, τότε είναι $1-1$. Άρα η εξίσωση $f(x)=0$ μπορεί να έχει το πολύ μία λύση.

\item \textbf{Λάθος}. Η ιδιότητα των ενδιάμεσων τιμών μεταξύ των τιμών στα άκρα δεν συνεπάγεται υποχρεωτικά συνέχεια.

\item \textbf{Σωστό}. Αν η $f$ είναι συνεχής σε διάστημα και δεν μηδενίζεται, τότε δεν μπορεί να αλλάζει πρόσημο. Διαφορετικά, από το Θεώρημα Bolzano θα είχε ρίζα.

\item \textbf{Λάθος}. Μια άρτια και συνεχής συνάρτηση μπορεί να έχει άπειρες ρίζες. Για παράδειγμα, η μηδενική συνάρτηση έχει άπειρες ρίζες.

\item \textbf{Λάθος}. Επειδή
\[
\left|\frac{\hm x}{x}\right|\leq \frac{1}{|x|},
\]
παίρνουμε
\[
\lim_{x\to-\infty}\frac{\hm x}{x}=0,
\]
όχι $-\infty$.

\item \textbf{Σωστό}. Αν $f(a)=f(\beta)$ και υπάρχει $\gamma\in(a,\beta)$ με $f(\gamma)\neq f(a)$, τότε, λόγω συνέχειας, κάποια ενδιάμεση τιμή λαμβάνεται δύο φορές. Άρα η $f$ δεν είναι $1-1$.

\item \textbf{Λάθος}. Για κάθε $x\in\mathbb{R}$ ισχύει
\[
x^4+3x^2+5>0.
\]
Άρα η εξίσωση δεν έχει πραγματική ρίζα.
\end{enumerate}

\subsection{Διαφορικός Λογισμός}

\begin{enumerate}[label=\textbf{2.\arabic*.}]
\item \textbf{Λάθος}. Η συνέχεια δεν συνεπάγεται παραγωγισιμότητα. Για παράδειγμα, η $f(x)=|x|$ είναι συνεχής στο $0$, αλλά δεν είναι παραγωγίσιμη εκεί.

\item \textbf{Λάθος}. Μια συνάρτηση μπορεί να είναι συνεχής αλλά όχι παραγωγίσιμη σε ένα σημείο. Για παράδειγμα, η $f(x)=|x|$ είναι συνεχής στο $0$, αλλά δεν είναι παραγωγίσιμη στο $0$.

\item \textbf{Σωστό}. Κάθε πολυωνυμική συνάρτηση είναι παραγωγίσιμη σε όλο το $\mathbb{R}$.

\item \textbf{Λάθος}. Μια συνάρτηση μπορεί να είναι παραγωγίσιμη χωρίς η παράγωγός της να είναι συνεχής.

\item \textbf{Λάθος}. Ο ρυθμός μεταβολής της ταχύτητας είναι η επιτάχυνση. Η παράγωγος της επιτάχυνσης είναι ο ρυθμός μεταβολής της επιτάχυνσης.

\item \textbf{Σωστό}. Αν η $f$ είναι παραγωγίσιμη στο $x_0$, τότε ο συντελεστής διεύθυνσης της εφαπτομένης της $C_f$ στο
\[
M(x_0,f(x_0))
\]
είναι ο αριθμός $f'(x_0)$.

\item \textbf{Λάθος}. Μια εφαπτομένη μπορεί να έχει και άλλα κοινά σημεία με την καμπύλη. Για παράδειγμα, η εφαπτομένη της $y=x^3$ στο $0$ είναι ο άξονας $x'x$, που τέμνει την καμπύλη στο σημείο επαφής αλλά η γενική απαγόρευση δεν ισχύει.

\item \textbf{Σωστό}. Είναι δυνατόν η ίδια ευθεία να εφάπτεται σε δύο διαφορετικά σημεία μιας γραφικής παράστασης, ανάλογα με τη μορφή της καμπύλης.

\item \textbf{Λάθος}. Η παράγωγος της $f(x)=5^x$ είναι
\[
f'(x)=5^x\ln5,
\]
όχι $x\cdot5^{x-1}$.

\item \textbf{Σωστό}. Για $f(x)=e^x$ έχουμε
\[
f(0)=1
\quad\text{και}\quad
f'(0)=1.
\]
Άρα η εφαπτομένη στο $A(0,1)$ είναι
\[
y-1=x,
\]
δηλαδή $y=x+1$.

\item \textbf{Λάθος}. Ο κανόνας γινομένου είναι
\[
(fg)'=f'g+fg',
\]
όχι $(fg)'=f'g'$.

\item \textbf{Λάθος}. Ισχύει
\[
\left(\frac{1}{f(x)}\right)'=-\frac{f'(x)}{(f(x))^2},
\]
όπου $f(x)\neq0$. Λείπει ο παράγοντας $f'(x)$.

\item \textbf{Λάθος}. Το $e^\pi$ είναι σταθερός αριθμός, άρα η παράγωγός του είναι $0$.

\item \textbf{Σωστό}. Παραγωγίζοντας τη σχέση
\[
(f(x))^2+(g(x))^2=1
\]
παίρνουμε
\[
2f(x)f'(x)+2g(x)g'(x)=0,
\]
οπότε
\[
f(x)f'(x)+g(x)g'(x)=0.
\]

\item \textbf{Λάθος}. Αν η $f$ είναι άρτια και παραγωγίσιμη στο $\mathbb{R}$, τότε η $f'$ είναι περιττή, όχι άρτια.

\item \textbf{Λάθος}. Αν η $f$ είναι περιττή, δεν είναι υποχρεωτικό να ισχύει $f'(0)=0$. Για παράδειγμα, η $f(x)=x$ είναι περιττή και $f'(0)=1$.

\item \textbf{Λάθος}. Μια παραγωγίσιμη και $1-1$ συνάρτηση μπορεί να έχει αντίστροφη που δεν είναι παραγωγίσιμη σε κάποιο σημείο. Για παράδειγμα, η $f(x)=x^3$ έχει αντίστροφη $\sqrt[3]{x}$, που δεν είναι παραγωγίσιμη στο $0$.

\item \textbf{Λάθος}. Το ότι δύο συναρτήσεις τέμνονται σε ένα σημείο σημαίνει μόνο ότι έχουν ίδια τιμή εκεί, όχι ότι έχουν και ίδια παράγωγο.

\item \textbf{Λάθος}. Αν η $C_f$ τέμνει τον άξονα $x'x$, τότε $f(x_0)=0$. Αυτό δεν σημαίνει ότι $f'(x_0)=0$. Για παράδειγμα, η $f(x)=x$ τέμνει τον άξονα στο $0$ και έχει εφαπτομένη $y=x$.

\item \textbf{Λάθος}. Κρίσιμα σημεία είναι εσωτερικά σημεία στα οποία η παράγωγος μηδενίζεται ή δεν υπάρχει. Άρα δεν είναι μόνο τα σημεία όπου $f'(x)=0$.

\item \textbf{Λάθος}. Στα άκρα ενός κλειστού διαστήματος αναζητούμε υποψήφιες θέσεις ολικών ακροτάτων, αλλά τα κρίσιμα σημεία είναι εσωτερικά σημεία.

\item \textbf{Λάθος}. Το Θεώρημα Fermat απαιτεί το ακρότατο να βρίσκεται σε εσωτερικό σημείο και η συνάρτηση να είναι παραγωγίσιμη εκεί. Για παράδειγμα, η $f(x)=|x|$ έχει τοπικό ελάχιστο στο $0$, αλλά δεν είναι παραγωγίσιμη στο $0$.

\item \textbf{Λάθος}. Η συνθήκη
\[
f'(x_0)=0
\]
δίνει πιθανή θέση ακροτάτου, όχι σίγουρη. Για παράδειγμα, η $f(x)=x^3$ έχει $f'(0)=0$, αλλά δεν έχει ακρότατο στο $0$.

\item \textbf{Σωστό}. Αν μια συνάρτηση παρουσιάζει τοπικό ακρότατο σε εσωτερικό σημείο $x_0$ και είναι παραγωγίσιμη εκεί, τότε από το Θεώρημα Fermat ισχύει
\[
f'(x_0)=0.
\]

\item \textbf{Σωστό}. Αν
\[
f'(x_0)=0,
\]
τότε η εφαπτομένη, εφόσον υπάρχει, έχει συντελεστή διεύθυνσης $0$. Άρα είναι παράλληλη στον άξονα $x'x$.

\item \textbf{Λάθος}. Το Θεώρημα Rolle εξασφαλίζει την ύπαρξη τουλάχιστον ενός σημείου με οριζόντια εφαπτομένη, όχι ακριβώς ενός.

\item \textbf{Λάθος}. Από τις υποθέσεις $f(a)=f(\beta)$ και παραγωγισιμότητα δεν προκύπτει ότι η $f$ είναι σταθερή. Για παράδειγμα, η $f(x)=x(1-x)$ στο $[0,1]$ έχει $f(0)=f(1)=0$, αλλά δεν είναι σταθερή.

\item \textbf{Σωστό}. Αν μια παραγωγίσιμη συνάρτηση έχει δύο διαφορετικές ρίζες $a,\beta$, τότε $f(a)=f(\beta)=0$. Από το Θεώρημα Rolle υπάρχει $\xi\in(a,\beta)$ με
\[
f'(\xi)=0.
\]

\item \textbf{Λάθος}. Το Θεώρημα Μέσης Τιμής εφαρμόζεται σε συναρτήσεις συνεχείς στο κλειστό διάστημα και παραγωγίσιμες στο ανοικτό, χωρίς να απαιτείται γνήσια μονοτονία.

\item \textbf{Σωστό}. Το Θεώρημα Rolle είναι ειδική περίπτωση του Θεωρήματος Μέσης Τιμής όταν
\[
f(a)=f(\beta).
\]

\item \textbf{Σωστό}. Το Θεώρημα Μέσης Τιμής εξασφαλίζει σημείο $\xi\in(a,\beta)$ τέτοιο ώστε
\[
f'(\xi)=\frac{f(\beta)-f(a)}{\beta-a}.
\]
Άρα η εφαπτομένη στο $\xi$ είναι παράλληλη στη χορδή των άκρων.

\item \textbf{Λάθος}. Αν το πεδίο ορισμού είναι ένωση ξένων διαστημάτων, η συνάρτηση μπορεί να είναι σταθερή σε κάθε διάστημα με διαφορετική σταθερά. Άρα δεν είναι υποχρεωτικά σταθερή σε όλο το πεδίο ορισμού.

\item \textbf{Λάθος}. Το $\mathbb{R}^*$ δεν είναι διάστημα αλλά ένωση δύο διαστημάτων. Μπορεί να υπάρχουν διαφορετικές σταθερές στα $(-\infty,0)$ και $(0,+\infty)$.

\item \textbf{Σωστό}. Αν
\[
f'(x)>0
\]
σε κάθε σημείο ενός ανοιχτού διαστήματος $\varDelta$, τότε από το Θεώρημα Μέσης Τιμής η $f$ είναι γνησίως αύξουσα στο $\varDelta$.

\item \textbf{Λάθος}. Μια γνησίως αύξουσα και παραγωγίσιμη συνάρτηση μπορεί να έχει μηδενική παράγωγο σε κάποιο σημείο. Για παράδειγμα, η $f(x)=x^3$ είναι γνησίως αύξουσα στο $\mathbb{R}$, αλλά $f'(0)=0$.

\item \textbf{Σωστό}. Υπάρχουν γνησίως αύξουσες συναρτήσεις των οποίων η παράγωγος μηδενίζεται σε άπειρα μεμονωμένα σημεία. Η μηδενική παράγωγος σε μεμονωμένα σημεία δεν αναιρεί κατ' ανάγκη τη γνήσια αύξηση.

\item \textbf{Λάθος}. Το ότι η εξίσωση $f'(x)=0$ έχει μόνο μία ρίζα δεν σημαίνει υποχρεωτικά ότι η $f$ έχει τοπικό ακρότατο. Για παράδειγμα, η $f(x)=x^3$ έχει μοναδική ρίζα της παραγώγου στο $0$, αλλά δεν έχει ακρότατο εκεί.

\item \textbf{Λάθος}. Ένα τοπικό ελάχιστο δεν είναι απαραίτητα μικρότερο από κάθε τοπικό μέγιστο. Τα τοπικά ακρότατα συγκρίνουν τιμές μόνο σε μια περιοχή του σημείου.

\item \textbf{Σωστό}. Η μέγιστη τιμή μιας συνεχούς συνάρτησης σε κλειστό διάστημα λαμβάνεται είτε σε άκρο είτε σε εσωτερικό σημείο. Αν λαμβάνεται σε εσωτερικό σημείο και υπάρχει παράγωγος, τότε το σημείο είναι κρίσιμο.

\item \textbf{Σωστό}. Μια συνάρτηση μπορεί να έχει άπειρα τοπικά ακρότατα. Για παράδειγμα, τριγωνομετρικές συναρτήσεις όπως η $\hm x$ έχουν άπειρα τοπικά μέγιστα και ελάχιστα στο $\mathbb{R}$.

\item \textbf{Λάθος}. Σημείο καμπής είναι σημείο στο οποίο αλλάζει η κυρτότητα της γραφικής παράστασης. Το ότι η εφαπτομένη διαπερνά την καμπύλη είναι γεωμετρική ένδειξη σε πολλές περιπτώσεις, όχι ο ορισμός.

\item \textbf{Λάθος}. Μπορεί να υπάρχει σημείο καμπής χωρίς να υπάρχει η $f''(x_0)$ ή χωρίς να ισχύει $f''(x_0)=0$. Η αλλαγή κυρτότητας είναι το ουσιαστικό κριτήριο.

\item \textbf{Λάθος}. Η συνθήκη
\[
f''(x_0)=0
\]
δεν αρκεί για σημείο καμπής. Χρειάζεται να αλλάζει κυρτότητα η $C_f$ εκατέρωθεν του $x_0$.

\item \textbf{Σωστό}. Αν
\[
f''(x)>0
\]
σε ένα διάστημα $\varDelta$, τότε η $f$ είναι κυρτή, δηλαδή στρέφει τα κοίλα προς τα άνω στο $\varDelta$.

\item \textbf{Σωστό}. Για κυρτή συνάρτηση, η γραφική παράσταση βρίσκεται πάνω από κάθε εφαπτομένη της, εκτός από το σημείο επαφής.

\item \textbf{Λάθος}. Το Θεώρημα Rolle εξασφαλίζει την ύπαρξη τουλάχιστον ενός $x_0\in(a,\beta)$ με
\[
f'(x_0)=0,
\]
όχι ακριβώς ενός.

\item \textbf{Λάθος}. Ο κανόνας De L'Hospital εφαρμόζεται άμεσα σε μορφές
\[
\frac{0}{0}
\quad\text{ή}\quad
\frac{\pm\infty}{\pm\infty},
\]
όχι απευθείας σε κάθε μορφή $0\cdot(\pm\infty)$ χωρίς μετασχηματισμό.

\item \textbf{Σωστό}. Όταν ισχύουν όλες οι προϋποθέσεις του κανόνα De L'Hospital και υπάρχει το όριο
\[
\lim \frac{f'(x)}{g'(x)},
\]
τότε
\[
\lim \frac{f(x)}{g(x)}=\lim \frac{f'(x)}{g'(x)}.
\]

\item \textbf{Λάθος}. Αν δεν υπάρχει το όριο
\[
\lim \frac{f'(x)}{g'(x)},
\]
δεν συμπεραίνουμε απαραίτητα ότι δεν υπάρχει το αρχικό όριο. Ο κανόνας απλώς δεν δίνει τότε αποτέλεσμα.

\item \textbf{Σωστό}. Κάθε πολυωνυμική συνάρτηση βαθμού τουλάχιστον $2$ δεν έχει οριζόντια ή πλάγια ασύμπτωτη, ούτε κατακόρυφη ασύμπτωτη, αφού είναι συνεχής σε όλο το $\mathbb{R}$.

\item \textbf{Λάθος}. Τα τοπικά μέγιστα δεν είναι πάντοτε μεγαλύτερα από τα τοπικά ελάχιστα. Τα τοπικά ακρότατα συγκρίνουν τιμές μόνο σε μια περιοχή του σημείου.

\item \textbf{Σωστό}. Αν
\[
f'(x)=g'(x)
\]
στο $\mathbb{R}$, τότε
\[
(f-g)'(x)=0
\]
στο $\mathbb{R}$. Άρα
\[
f(x)-g(x)=c,
\]
οπότε οι γραφικές παραστάσεις προκύπτουν με κατακόρυφη μετατόπιση η μία από την άλλη.

\item \textbf{Λάθος}. Η ύπαρξη δύο σημείων καμπής δεν συνεπάγεται υποχρεωτικά την ύπαρξη τοπικού ακροτάτου ανάμεσά τους.

\item \textbf{Σωστό}. Για παράδειγμα, η συνάρτηση
\[
f(x)=e^{-x}-x
\]
είναι κυρτή στο $\mathbb{R}$, αφού
\[
f''(x)=e^{-x}>0,
\]
και έχει πλάγια ασύμπτωτη την ευθεία $y=-x$ στο $+\infty$.

\item \textbf{Λάθος}. Η συνάρτηση $e^x$ έχει
\[
\lim_{x\to-\infty}e^x=0,
\]
αλλά δεν παίρνει ποτέ την τιμή $0$. Άρα το $0$ είναι κάτω φράγμα, όχι ελάχιστο.

\item \textbf{Σωστό}. Αν $f(x)=c$ για κάθε $x\in\mathbb{R}$, τότε
\[
f'(x)=0
\]
για κάθε $x\in\mathbb{R}$.

\item \textbf{Λάθος}. Μια γνησίως αύξουσα και παραγωγίσιμη συνάρτηση μπορεί να έχει μηδενική παράγωγο σε σημείο. Για παράδειγμα, η $f(x)=x^3$ είναι γνησίως αύξουσα, αλλά $f'(0)=0$.

\item \textbf{Σωστό}. Το Θεώρημα Fermat αφορά τοπικά ακρότατα που παρουσιάζονται σε εσωτερικά σημεία του πεδίου ορισμού, όπου η συνάρτηση είναι παραγωγίσιμη.

\item \textbf{Λάθος}. Αν μια συνάρτηση είναι συνεχής σε όλο το $\mathbb{R}$, τότε δεν μπορεί να έχει κατακόρυφη ασύμπτωτη, αφού δεν μπορεί να έχει άπειρο όριο σε πραγματικό σημείο του πεδίου της.

\item \textbf{Σωστό}. Για $x\neq0$ ισχύει
\[
(\ln|x|)'=\frac{1}{x}.
\]

\item \textbf{Λάθος}. Το Θεώρημα Bolzano ανήκει στα θεωρήματα συνέχειας και δεν είναι άμεση συνέπεια του Θεωρήματος Μέσης Τιμής.

\item \textbf{Λάθος}. Η συνθήκη
\[
f''(x)>0
\]
σημαίνει ότι η $f$ είναι κυρτή, όχι ότι είναι γνησίως αύξουσα. Για παράδειγμα, η $f(x)=x^2$ έχει $f''(x)=2>0$, αλλά δεν είναι γνησίως αύξουσα σε όλο το $\mathbb{R}$.

\item \textbf{Σωστό}. Για $f(x)=\hm x$ έχουμε
\[
f(0)=0
\quad\text{και}\quad
f'(0)=\syn0=1.
\]
Άρα η εφαπτομένη στο $x=0$ είναι
\[
y=x.
\]

\item \textbf{Σωστό}. Αν η εφαπτομένη της $C_f$ στο $x_0$ είναι
\[
y=2x-1,
\]
τότε ο συντελεστής διεύθυνσής της είναι $2$, άρα $f'(x_0)=2$, και επειδή περνά από το σημείο επαφής,
\[
f(x_0)=2x_0-1.
\]

\item \textbf{Σωστό}. Σε σημείο καμπής η γραφική παράσταση αλλάζει κυρτότητα. Όταν υπάρχει εφαπτομένη στο σημείο αυτό, η εφαπτομένη διαπερνά την καμπύλη.

\item \textbf{Λάθος}. Μια γνησίως αύξουσα και παραγωγίσιμη συνάρτηση μπορεί να έχει μηδενική παράγωγο σε σημείο. Για παράδειγμα, η $f(x)=x^3$ είναι γνησίως αύξουσα, αλλά
\[
f'(0)=0.
\]

\item \textbf{Σωστό}. Αν η $f$ είναι περιττή και παραγωγίσιμη, τότε η $f'$ είναι άρτια. Πράγματι, από
\[
f(-x)=-f(x)
\]
παραγωγίζοντας παίρνουμε
\[
f'(-x)=f'(x).
\]

\item \textbf{Σωστό}. Για $f(x)=x^3$ έχουμε
\[
f'(x)=3x^2,
\]
που είναι άρτια συνάρτηση.

\item \textbf{Λάθος}. Μια πολυωνυμική συνάρτηση 3ου βαθμού δεν έχει υποχρεωτικά τοπικό μέγιστο και τοπικό ελάχιστο. Για παράδειγμα, η $f(x)=x^3$ δεν έχει τοπικά ακρότατα.

\item \textbf{Σωστό}. Μια πολυωνυμική συνάρτηση περιττού βαθμού τουλάχιστον $3$ έχει δεύτερη παράγωγο πολυώνυμο περιττού βαθμού, άρα έχει τουλάχιστον μία πραγματική ρίζα και παρουσιάζει τουλάχιστον ένα σημείο καμπής.

\item \textbf{Σωστό}. Αν
\[
\lim_{x\to x_0}\frac{f(x)-f(x_0)}{x-x_0}=-\infty,
\]
τότε το όριο του λόγου μεταβολής δεν είναι πραγματικός αριθμός. Άρα η $f$ δεν είναι παραγωγίσιμη στο $x_0$.

\item \textbf{Σωστό}. Η παράγωγος $f'$ ορίζεται μόνο στα σημεία του πεδίου ορισμού της $f$ στα οποία υπάρχει το όριο του λόγου μεταβολής. Άρα
\[
D_{f'}\subseteq D_f.
\]

\item \textbf{Σωστό}. Αν η $f$ είναι παραγωγίσιμη στο $x_0$, τότε υπάρχει μοναδικός συντελεστής διεύθυνσης $f'(x_0)$, άρα υπάρχει μοναδική εφαπτομένη στο σημείο αυτό.

\item \textbf{Λάθος}. Η μορφή
\[
\frac{\pm\infty}{\pm\infty}
\]
μπορεί να αντιμετωπιστεί με De L'Hospital μόνο όταν ισχύουν οι προϋποθέσεις του κανόνα. Δεν εφαρμόζεται μηχανικά χωρίς έλεγχο.

\item \textbf{Σωστό}. Αν η $f''$ αλλάζει πρόσημο εκατέρωθεν του $x_0$ και υπάρχει εφαπτομένη της $C_f$ στο σημείο αυτό, τότε αλλάζει η κυρτότητα της $f$. Άρα το $x_0$ είναι θέση σημείου καμπής.

\item \textbf{Λάθος}. Η μη ύπαρξη της $f''(x_0)$ δεν αποκλείει τοπικό ακρότατο. Για παράδειγμα, η $f(x)=|x|$ έχει τοπικό ελάχιστο στο $0$, αλλά δεν είναι δύο φορές παραγωγίσιμη εκεί.

\item \textbf{Λάθος}. Το Θεώρημα Fermat εφαρμόζεται μόνο όταν η συνάρτηση είναι παραγωγίσιμη στο εσωτερικό σημείο όπου παρουσιάζει τοπικό ακρότατο.

\item \textbf{Σωστό}. Η συνάρτηση $f(x)=x^3$ είναι γνησίως αύξουσα στο $\mathbb{R}$, παρόλο που
\[
f'(0)=0.
\]

\item \textbf{Λάθος}. Αν το πεδίο ορισμού είναι ένωση διαστημάτων, το $f'(x)>0$ εξασφαλίζει γνήσια αύξηση σε κάθε διάστημα χωριστά, όχι απαραίτητα σε όλο το πεδίο ορισμού.

\item \textbf{Σωστό}. Σε θέση τοπικού ελάχιστου είναι δυνατόν η συνάρτηση να μην είναι παραγωγίσιμη. Κλασικό παράδειγμα είναι η $f(x)=|x|$ στο $x=0$.

\item \textbf{Λάθος}. Μια κυρτή συνάρτηση στο $\mathbb{R}$ δεν είναι υποχρεωτικό να τείνει στο $+\infty$ όταν $x\to+\infty$. Για παράδειγμα, η $f(x)=e^{-x}$ είναι κυρτή και
\[
\lim_{x\to+\infty}e^{-x}=0.
\]

\item \textbf{Λάθος}. Η απόσταση δύο κινητών σε παράλληλους άξονες δεν ακροτατοποιείται γενικά μόνο όταν εξισώνονται οι ταχύτητές τους. Χρειάζεται να παραγωγιστεί η ίδια η απόσταση ή το τετράγωνό της.

\item \textbf{Λάθος}. Ο ρυθμός μεταβολής της τετμημένης είναι $x'(t)$, αλλά δεν είναι πάντα ίσος με $1$. Εξαρτάται από τον τρόπο κίνησης.

\item \textbf{Σωστό}. Σε ευθύγραμμη κίνηση, αν η ταχύτητα και η επιτάχυνση είναι ετερόσημες, τότε το μέτρο της ταχύτητας μειώνεται. Άρα το κινητό επιβραδύνεται.

\item \textbf{Λάθος}. Το Θεώρημα Rolle δίνει σχέση από τις ρίζες της $f$ προς ρίζες της $f'$, όχι αντίστροφα. Οι ρίζες της $f'$ δεν διαχωρίζονται υποχρεωτικά από ρίζες της $f$.

\item \textbf{Λάθος}. Μεταξύ δύο διαδοχικών ριζών μιας παραγωγίσιμης συνάρτησης υπάρχει τουλάχιστον μία ρίζα της παραγώγου, όχι ακριβώς μία.

\item \textbf{Λάθος}. Αν $f'(x_0)=g'(x_0)=0$, τότε οι δύο γραφικές παραστάσεις έχουν οριζόντια εφαπτομένη στο κοινό σημείο, αλλά δεν σημαίνει ότι το σημείο είναι ακρότατο και για τις δύο συναρτήσεις.

\item \textbf{Λάθος}. Για $f(x)=e^{-x}$ έχουμε
\[
\lim_{x\to-\infty}e^{-x}=+\infty.
\]
Άρα δεν υπάρχει οριζόντια ασύμπτωτη στο $-\infty$.

\item \textbf{Λάθος}. Ο κανόνας De L'Hospital εφαρμόζεται και σε απροσδιοριστίες της μορφής
\[
\frac{\pm\infty}{\pm\infty},
\]
όχι αποκλειστικά στη μορφή $\frac00$.

\item \textbf{Σωστό}. Μια γραφική παράσταση μπορεί να τέμνει την πλάγια ασύμπτωτή της σε άπειρα σημεία. Για παράδειγμα, η $f(x)=x+\frac{\hm x}{x}$ έχει πλάγια ασύμπτωτη $y=x$ στο $+\infty$ και την τέμνει άπειρες φορές.

\item \textbf{Σωστό}. Αν υπάρχει πλάγια ασύμπτωτη στο $+\infty$ με μη μηδενική κλίση, τότε δεν μπορεί ταυτόχρονα να υπάρχει οριζόντια ασύμπτωτη στο $+\infty$.

\item \textbf{Λάθος}. Το όριο
\[
\lim_{x\to+\infty}\frac{f(x)}{x}=\lambda
\]
δεν αρκεί για ύπαρξη πλάγιας ή οριζόντιας ασύμπτωτης. Χρειάζεται επιπλέον να υπάρχει πεπερασμένο το όριο
\[
\lim_{x\to+\infty}\bigl(f(x)-\lambda x\bigr).
\]

\item \textbf{Σωστό}. Για κύκλο με εμβαδόν
\[
E=\pi r^2
\]
έχουμε
\[
\frac{\d E}{\d r}=2\pi r,
\]
που είναι το μήκος της περιφέρειας.

\item \textbf{Σωστό}. Αν η $f'$ είναι γνησίως αύξουσα σε ένα διάστημα, τότε η $f$ είναι κυρτή στο διάστημα αυτό.

\item \textbf{Σωστό}. Υπό τις προϋποθέσεις του Θεωρήματος Rolle, υπάρχει τουλάχιστον ένα $x_0\in(a,\beta)$ με
\[
f'(x_0)=0.
\]
Άρα η $C_f$ δέχεται οριζόντια εφαπτομένη στο σημείο αυτό.

\item \textbf{Λάθος}. Από τις προϋποθέσεις συνέχειας και παραγωγισιμότητας δεν εξασφαλίζεται υποχρεωτικά οριζόντια εφαπτομένη. Θα χρειαζόταν, για παράδειγμα, η επιπλέον συνθήκη $f(a)=f(\beta)$, όπως στο Θεώρημα Rolle.

\item \textbf{Λάθος}. Μια συνάρτηση μπορεί να μην είναι συνεχής σε σημείο και να έχει εκεί τοπικό ακρότατο. Για παράδειγμα, αν $f(0)=1$ και $f(x)=0$ για $x\neq0$, τότε η $f$ έχει τοπικό μέγιστο στο $0$ αλλά δεν είναι συνεχής εκεί.

\item \textbf{Σωστό}. Αν η $f'$ παρουσιάζει τοπικό μέγιστο στο $x_0$, τότε η $f'$ μπορεί να αλλάζει μονοτονία εκεί. Αυτό κάνει το $x_0$ πιθανή θέση αλλαγής κυρτότητας της $f$, δηλαδή πιθανή θέση σημείου καμπής.

\item \textbf{Σωστό}. Για κυρτή συνάρτηση, η γραφική παράσταση βρίσκεται πάνω από κάθε εφαπτομένη της. Άρα η τεταγμένη της καμπύλης είναι μεγαλύτερη ή ίση από την αντίστοιχη τεταγμένη της εφαπτομένης.

\item \textbf{Λάθος}. Αν δεν υπάρχει συνέχεια, μια συνάρτηση χωρίς ρίζες μπορεί να παίρνει και θετικές και αρνητικές τιμές. Για παράδειγμα,
\[
f(x)=
\begin{cases}
1,&x\in\mathbb{Q},\\
-1,&x\notin\mathbb{Q}.
\end{cases}
\]
Δεν έχει ρίζες, αλλά δεν διατηρεί σταθερό πρόσημο.
\end{enumerate}

\subsection{Ολοκληρωτικός Λογισμός}

\begin{enumerate}[label=\textbf{3.\arabic*.}]
\item \textbf{Λάθος}. Χωρίς υπόθεση συνέχειας, μπορεί να ισχύει $f(x)\geq0$ και το ολοκλήρωμα να είναι $0$, ενώ η $f$ να μη μηδενίζεται σε μεμονωμένα σημεία.

\item \textbf{Λάθος}. Αν η $f$ παίρνει και αρνητικές τιμές, το ολοκλήρωμα δίνει προσημασμένο εμβαδόν. Το γεωμετρικό εμβαδόν δίνεται από
\[
\dint_a^\beta |f(x)|\,\d x.
\]

\item \textbf{Λάθος}. Ίσα ολοκληρώματα δεν σημαίνουν υποχρεωτικά ίσα γεωμετρικά εμβαδά, γιατί μπορεί να υπάρχουν θετικά και αρνητικά τμήματα που αλληλοαναιρούνται.

\item \textbf{Λάθος}. Το
\[
\dint_{-1}^{1}x^3\,\d x=0
\]
είναι προσημασμένο ολοκλήρωμα. Το εμβαδόν δίνεται από
\[
\dint_{-1}^{1}|x^3|\,\d x.
\]

\item \textbf{Σωστό}. Αν $a<\beta$ και η συνεχής $f$ δεν είχε ρίζα στο $(a,\beta)$, τότε θα διατηρούσε σταθερό πρόσημο στο εσωτερικό και το ολοκλήρωμα δεν θα μπορούσε να είναι $0$.

\item \textbf{Σωστό}. Κάθε συνεχής συνάρτηση σε κλειστό διάστημα $[a,\beta]$ είναι ολοκληρώσιμη σε αυτό.

\item \textbf{Σωστό}. Είναι ο τύπος της παραγοντικής ολοκλήρωσης:
\[
\dint_a^\beta f(x)g'(x)\,\d x
=f(\beta)g(\beta)-f(a)g(a)-\dint_a^\beta f'(x)g(x)\,\d x.
\]

\item \textbf{Σωστό}. Αν $f(x)<0$ στο $[a,\beta]$, το εμβαδόν είναι
\[
-\dint_a^\beta f(x)\,\d x
=
\dint_\beta^a f(x)\,\d x.
\]
Άρα πράγματι δίνεται από το ολοκλήρωμα με αντεστραμμένα άκρα.

\item \textbf{Λάθος}. Από το
\[
\dint_a^\beta f(x)\,\d x>0
\]
δεν προκύπτει ότι $f(x)>0$ για κάθε $x\in[a,\beta]$. Η $f$ μπορεί να παίρνει και αρνητικές τιμές, αρκεί η συνολική προσημασμένη επιφάνεια να είναι θετική.

\item \textbf{Λάθος}. Η ισότητα των ολοκληρωμάτων δεν συνεπάγεται ισότητα των συναρτήσεων. Διαφορετικές συναρτήσεις μπορεί να έχουν το ίδιο ολοκλήρωμα στο ίδιο διάστημα.

\item \textbf{Λάθος}. Γενικά δεν ισχύει
\[
\dint_a^\beta f(x)g(x)\,\d x
=
\left(\dint_a^\beta f(x)\,\d x\right)
\left(\dint_a^\beta g(x)\,\d x\right).
\]
Το ολοκλήρωμα δεν μοιράζεται στο γινόμενο με αυτόν τον τρόπο.

\item \textbf{Σωστό}. Αν $f(x)\geq g(x)$ στο $[a,\beta]$, τότε η κατακόρυφη απόσταση των δύο καμπυλών είναι
\[
f(x)-g(x)\geq0.
\]
Άρα το εμβαδόν μεταξύ τους είναι
\[
\dint_a^\beta (f(x)-g(x))\,\d x,
\]
ανεξάρτητα από τη θέση τους ως προς τον άξονα $x'x$.

\item \textbf{Λάθος}. Η ιδιότητα
\[
\dint_a^\beta c f(x)\,\d x
=c\dint_a^\beta f(x)\,\d x
\]
ισχύει για κάθε σταθερά $c\in\mathbb{R}$, όχι μόνο για $c>0$.

\item \textbf{Σωστό}. Για κάθε ολοκληρώσιμη συνάρτηση ισχύει
\[
\dint_a^a f(x)\,\d x=0.
\]

\item \textbf{Λάθος}. Αν η $f$ είναι άρτια, τότε
\[
\dint_{-a}^{a} f(x)\,\d x=2\dint_0^a f(x)\,\d x,
\]
όχι υποχρεωτικά $0$.

\item \textbf{Σωστό}. Αν η $f$ είναι περιττή και συνεχής, τότε τα προσημασμένα εμβαδά στο $[-a,0]$ και στο $[0,a]$ αλληλοαναιρούνται. Άρα
\[
\dint_{-a}^{a} f(x)\,\d x=0.
\]

\item \textbf{Λάθος}. Το εμβαδόν επίπεδου χωρίου είναι πάντοτε μη αρνητικός αριθμός. Αρνητικό μπορεί να είναι το προσημασμένο ολοκλήρωμα, όχι το εμβαδόν.

\item \textbf{Λάθος}. Από το
\[
\dint_a^\beta f(x)\,\d x\geq0
\]
δεν αποκλείεται η $f$ να παίρνει αρνητικές τιμές. Μπορεί οι θετικές περιοχές να υπερισχύουν.

\item \textbf{Σωστό}. Αν $f(x)\geq0$ στο $[a,\beta]$ και $a<\beta$, τότε από τη μονοτονία του ολοκληρώματος
\[
\dint_a^\beta f(x)\,\d x\geq0.
\]

\item \textbf{Σωστό}. Η τιμή του ορισμένου ολοκληρώματος εξαρτάται από τα άκρα ολοκλήρωσης $a,\beta$, καθώς αυτά καθορίζουν το διάστημα ολοκλήρωσης.

\item \textbf{Λάθος}. Όπως είναι γραμμένη, η σχέση με άκρο $b$ δεν ισχύει γενικά. Η σωστή ιδιότητα είναι
\[
\dint_a^\beta f(x)\,\d x
=
-\dint_\beta^a f(x)\,\d x.
\]

\item \textbf{Λάθος}. Η αθροιστική ιδιότητα
\[
\dint_a^\beta f(x)\,\d x+\dint_\beta^\gamma f(x)\,\d x
=\dint_a^\gamma f(x)\,\d x
\]
ισχύει γενικά για προσανατολισμένα ολοκληρώματα, όχι μόνο όταν $a<\beta<\gamma$.

\item \textbf{Σωστό}. Αν
\[
\dint_a^\beta |f(x)|\,\d x=0
\]
και η $f$ είναι συνεχής, τότε $|f(x)|=0$ σε όλο το $[a,\beta]$. Άρα
\[
f(x)=0
\]
για κάθε $x\in[a,\beta]$.

\item \textbf{Σωστό}. Από το Θεμελιώδες Θεώρημα του Ολοκληρωτικού Λογισμού,
\[
\dint_a^\beta f'(x)\,\d x=f(\beta)-f(a),
\]
εφόσον η $f'$ είναι συνεχής στο $[a,\beta]$.

\item \textbf{Σωστό}. Η μεταβλητή ολοκλήρωσης είναι βουβή. Άρα
\[
\dint_0^1 f(x)\,\d x
\quad\text{και}\quad
\dint_0^1 f(t)\,\d t
\]
παριστάνουν τον ίδιο πραγματικό αριθμό.

\item \textbf{Λάθος}. Αν η $f'$ είναι συνεχής, τότε
\[
\dint_a^\beta f'(x)\,\d x=f(\beta)-f(a).
\]
Αυτό είναι προσημασμένη μεταβολή, όχι πάντα η απόσταση των τιμών, η οποία είναι
\[
|f(\beta)-f(a)|.
\]

\item \textbf{Σωστό}. Πράγματι,
\[
\left(\frac{e^{2x}}{2}\right)'=e^{2x}.
\]
Άρα η $\frac{e^{2x}}{2}$ είναι αρχική της $e^{2x}$.

\item \textbf{Λάθος}. Η σωστή σχέση είναι
\[
\dint_a^\beta f'(x)\,\d x=f(\beta)-f(a),
\]
όχι $f(a)-f(\beta)$.

\item \textbf{Σωστό}. Αν $c>0$ και $a<\beta$, τότε
\[
\dint_a^\beta c\,\d x=c(\beta-a),
\]
που είναι το εμβαδόν ορθογωνίου με ύψος $c$ και βάση $\beta-a$.

\item \textbf{Σωστό}. Το γεωμετρικό εμβαδόν είναι πάντοτε μη αρνητικός αριθμός, ανεξάρτητα από τη θέση των καμπυλών ως προς τον άξονα $x'x$.

\item \textbf{Σωστό}. Αν $f(x)\neq0$ στο $[a,\beta]$ και η $f'$ είναι συνεχής, τότε
\[
\left(\ln|f(x)|\right)'=\frac{f'(x)}{f(x)}.
\]
Άρα
\[
\dint_a^\beta \frac{f'(x)}{f(x)}\,\d x
=\ln|f(\beta)|-\ln|f(a)|.
\]

\item \textbf{Σωστό}. Από τον κανόνα γινομένου,
\[
(f(x)g(x))'=f'(x)g(x)+f(x)g'(x).
\]
Ολοκληρώνοντας στο $[a,\beta]$, παίρνουμε
\[
\dint_a^\beta f'(x)g(x)\,\d x+\dint_a^\beta f(x)g'(x)\,\d x
=f(\beta)g(\beta)-f(a)g(a).
\]

\item \textbf{Σωστό}. Έχουμε
\[
\dint_1^e \frac{1}{x}\,\d x
=\left[\ln x\right]_1^e
=1.
\]

\item \textbf{Σωστό}. Αν $f(x)\geq g(x)$, τότε η κατακόρυφη απόσταση των καμπυλών είναι
\[
f(x)-g(x)\geq0.
\]
Άρα το εμβαδόν δίνεται από
\[
\dint_a^\beta (f(x)-g(x))\,\d x,
\]
ακόμα κι αν οι δύο καμπύλες βρίσκονται κάτω από τον άξονα $x'x$.

\item \textbf{Σωστό}. Αν η συνεχής $f$ δεν τέμνει τον άξονα $x'x$, τότε διατηρεί σταθερό πρόσημο. Αν $f>0$, το εμβαδόν είναι
\[
\dint_a^\beta f(x)\,\d x,
\]
ενώ αν $f<0$, είναι
\[
-\dint_a^\beta f(x)\,\d x.
\]

\item \textbf{Σωστό}. Με την αντικατάσταση
\[
u=g(x)
\]
και $\d u=g'(x)\,\d x$, τα άκρα γίνονται $g(a)$ και $g(\beta)$. Άρα
\[
\dint_a^\beta f(g(x))g'(x)\,\d x
=
\dint_{g(a)}^{g(\beta)} f(u)\,\d u.
\]

\item \textbf{Λάθος}. Το ότι η $f$ είναι $1-1$ και παραγωγίσιμη δεν σημαίνει ότι το εμβαδόν που σχηματίζεται με τις $C_f$, $C_{f^{-1}}$ και την ευθεία $y=x$ είναι πάντοτε $0$.

\item \textbf{Λάθος}. Ισχύει
\[
(f(x)+g(x))^2=f^2(x)+2f(x)g(x)+g^2(x).
\]
Άρα γενικά λείπει ο όρος
\[
2\dint_a^\beta f(x)g(x)\,\d x.
\]

\item \textbf{Σωστό}. Αν δεν δίνονται κατακόρυφα όρια, για να βρούμε το κλειστό χωρίο μεταξύ των $C_f$ και $C_g$, πρέπει να βρούμε τα σημεία τομής τους.

\item \textbf{Σωστό}. Το ολοκλήρωμα
\[
E=\dint_a^\beta |f(x)|\,\d x
\]
υπολογίζει το συνολικό εμβαδόν μεταξύ της $C_f$ και του άξονα $x'x$, ακόμη κι αν η $f$ αλλάζει πρόσημο.

\item \textbf{Σωστό}. Έχουμε
\[
\dint_a^\beta e^{-x}\,\d x
=\left[-e^{-x}\right]_a^\beta
=e^{-a}-e^{-\beta}.
\]

\item \textbf{Σωστό}. Η συνάρτηση
\[
x^5+\hm x
\]
είναι περιττή. Άρα το ολοκλήρωμά της σε συμμετρικό διάστημα $[-a,a]$ ισούται με $0$.

\item \textbf{Λάθος}. Με την αλλαγή μεταβλητής $u=1-x$ παίρνουμε
\[
\dint_0^1 f(1-x)\,\d x=\dint_0^1 f(u)\,\d u.
\]
Άρα το αποτέλεσμα είναι $1$, όχι $-1$.

\item \textbf{Λάθος}. Η μεταβλητή ολοκλήρωσης είναι βουβή. Η τιμή του ολοκληρώματος δεν εξαρτάται από το αν γράφουμε $x$, $t$ ή άλλο γράμμα.

\item \textbf{Σωστό}. Αν $0<a<\beta$, τότε
\[
\frac{1}{x^2}>0
\]
στο $[a,\beta]$. Άρα
\[
\dint_a^\beta \frac{1}{x^2}\,\d x>0.
\]

\item \textbf{Σωστό}. Για $a<\beta$,
\[
\dint_a^\beta 1\,\d x=\beta-a,
\]
που είναι η απόσταση των σημείων $a$ και $\beta$ στον άξονα.

\item \textbf{Λάθος}. Το ολοκλήρωμα
\[
\dint_a^\beta f'(x)\,\d x
\]
δίνει προσημασμένη μεταβολή $f(\beta)-f(a)$, όχι πάντα γεωμετρικό εμβαδόν.

\item \textbf{Σωστό}. Αν $f$ είναι συνεχής και $f(x)\geq0$ στο $[a,\beta]$, τότε
\[
\dint_a^\beta f(x)\,\d x
\]
ταυτίζεται με το εμβαδόν του χωρίου κάτω από την $C_f$ και πάνω από τον άξονα $x'x$.

\item \textbf{Σωστό}. Από τη μονοτονία του ολοκληρώματος, αν
\[
m\leq f(x)\leq M
\]
στο $[a,\beta]$, τότε
\[
m(\beta-a)\leq\dint_a^\beta f(x)\,\d x\leq M(\beta-a).
\]

\item \textbf{Λάθος}. Το χωρίο μπορεί να μην περιγράφεται πάντα από το ολοκλήρωμα
\[
\dint_0^{x_0} f(x)\,\d x.
\]
Χρειάζεται έλεγχος προσήμου και σαφής περιγραφή των ορίων του χωρίου.

\item \textbf{Σωστό}. Έχουμε
\[
\dint_1^2 \frac{2x}{x^2+1}\,\d x
=\left[\ln(x^2+1)\right]_1^2
=\ln5-\ln2.
\]

\item \textbf{Λάθος}. Αν η $f$ είναι περιττή, το προσημασμένο ολοκλήρωμα στο $[-a,a]$ είναι $0$, αλλά το εμβαδόν είναι
\[
\dint_{-a}^a |f(x)|\,\d x,
\]
που δεν είναι υποχρεωτικά $0$.

\item \textbf{Σωστό}. Στο $[0,1]$ ισχύει $x\geq x^2$, άρα το εμβαδόν είναι
\[
\dint_0^1 (x-x^2)\,\d x=\frac16,
\]
θετικός αριθμός μικρότερος της μονάδας.

\item \textbf{Σωστό}. Αν
\[
\dint_1^3 f(x)\,\d x>0,
\]
τότε δεν γίνεται να ισχύει $f(x)\leq0$ για όλα τα $x\in[1,3]$. Άρα υπάρχει τουλάχιστον ένα $x_0$ με $f(x_0)>0$.

\item \textbf{Σωστό}. Με αντικατάσταση $u=f(x)$ έχουμε $\d u=f'(x)\,\d x$, οπότε
\[
\dint_a^\beta e^{f(x)}f'(x)\,\d x
=e^{f(\beta)}-e^{f(a)}.
\]

\item \textbf{Σωστό}. Η καμπύλη $y=\ln x$ τέμνει τον άξονα $x'x$ στο $x=1$. Άρα το χωρίο με τον άξονα $x'x$ και την ευθεία $x=e$ έχει εμβαδόν
\[
\dint_1^e \ln x\,\d x.
\]

\item \textbf{Σωστό}. Έχουμε
\[
\dint_0^\pi \syn x\,\d x
=\left[\hm x\right]_0^\pi
=0-0=0.
\]

\item \textbf{Σωστό}. Με αντικατάσταση $u=\ln x$ έχουμε $\d u=\frac1x\,\d x$. Άρα
\[
\dint_1^e \frac{\ln x}{x}\,\d x
=\dint_0^1 u\,\d u
=\frac12.
\]

\item \textbf{Λάθος}. Ίσα ορισμένα ολοκληρώματα δεν συνεπάγονται ότι οι συναρτήσεις είναι ίσες. Διαφορετικές συναρτήσεις μπορεί να έχουν την ίδια προσημασμένη επιφάνεια.

\item \textbf{Σωστό}. Στο διάστημα $[0,2\pi]$ τα εμβαδά της $\hm x$ πάνω και κάτω από τον άξονα $x'x$ είναι ίσα και αντίθετα ως προσημασμένα ολοκληρώματα. Άρα
\[
\dint_0^{2\pi}\hm x\,\d x=0.
\]

\item \textbf{Σωστό}. Αν $f(x)=c$ με $c>0$, τότε στο $[a,\beta]$ το αντίστοιχο χωρίο είναι ορθογώνιο με ύψος $c$ και βάση $\beta-a$.

\item \textbf{Λάθος}. Στο ορισμένο ολοκλήρωμα, όταν κάνουμε αντικατάσταση, μπορούμε είτε να αλλάξουμε τα όρια είτε να επιστρέψουμε στην αρχική μεταβλητή πριν αντικαταστήσουμε τα αρχικά όρια. Δεν είναι υποχρεωτική πάντοτε η αλλαγή των ορίων.

\item \textbf{Σωστό}. Από την ιδιότητα αλλαγής άκρων,
\[
\dint_\beta^a f(x)\,\d x=-\dint_a^\beta f(x)\,\d x.
\]
Άρα
\[
\dint_a^\beta f(x)\,\d x+\dint_\beta^a f(x)\,\d x=0.
\]

\item \textbf{Σωστό}. Με την αντικατάσταση $u=1-x$ παίρνουμε
\[
\dint_0^1 f(1-x)\,\d x=\dint_0^1 f(u)\,\d u.
\]
Άρα
\[
\dint_0^1 f(x)\,\d x=\dint_0^1 f(1-x)\,\d x.
\]

\item \textbf{Λάθος}. Το εμβαδόν μεταξύ δύο συναρτήσεων δίνεται γενικά από
\[
\dint_a^\beta |f(x)-g(x)|\,\d x,
\]
ή, αν ξέρουμε ποια είναι πάνω, από το ολοκλήρωμα της διαφοράς τους. Δεν δίνεται απαραίτητα από
\[
\dint_a^\beta \bigl(|f(x)|-|g(x)|\bigr)\,\d x.
\]
\end{enumerate}

%=========== ΤΕΛΟΣ ΑΠΑΝΤΗΣΕΩΝ 1.1-1.85, 2.1-2.100 ΚΑΙ 3.1-3.65 ===========
